\section{Introduction}

Estimating the geographic origin of a photograph from its pixels alone has
progressed from retrieval over geotagged collections~\cite{hays2008im2gps},
through classification over discretised geographic
cells~\cite{weyand2016planet,pramanick2022translocator}, to contrastive
alignment between images and continuous GPS coordinates.
GeoCLIP~\cite{cepeda2023geoclip} exemplifies the latter paradigm: a
CLIP~\cite{radford2021clip} vision tower is aligned with a location encoder,
and inference reduces to nearest-neighbour retrieval over a gallery of
$10^{5}$ GPS coordinates. Such models attain accuracies that render them
practically useful and, simultaneously, consequential, as the same
capability supports both the verification of news imagery and the inference
of location from photographs shared without geotags.

Despite this, the evidential basis of an individual prediction remains
largely opaque. It is widely assumed that these models exploit road
markings, vegetation, licence plates and shopfront text, but the supporting
evidence is predominantly anecdotal or derived from attention maps.
Attention distributions are correlational, and have been shown to be an
unreliable indication of the quantities upon which a model
depends~\cite{jain2019attention}; the same limitation applies to
gradient-based saliency~\cite{selvaraju2017gradcam} and to attention
rollout~\cite{abnar2020rollout}. The question of practical interest ---
whether a prediction would persist if the model were unable to attend to a
particular region --- is causal in nature, and therefore requires
intervention upon the model rather than inspection of it. Such analysis is
established practice in the study of language models, where ablating
specific attention pathways identifies the internal routes along which
factual information propagates~\cite{meng2022rome,geva2023dissecting}. The
present work transfers that methodology to image geolocation and renders it
interactive, so that the effect of removing a given visual cue may be
observed directly.

We present GeoProbe, whose contributions are threefold. First, an
interactive system applying region-conditioned, per-layer attention
intervention to the GeoCLIP vision tower at inference time, requiring
neither training, gradients nor architectural modification, and presenting
the unmodified and intervened predictions jointly on a map. Second, an
object-proposal workflow in which WeDetect-Uni~\cite{fu2026wedetect}
supplies category-agnostic foreground regions without requiring a prompt for
each image, enabling the same intervention to be evaluated at dataset scale.
Third, a quantitative study on Img2GPS3k~\cite{vo2017revisiting} that uses a
discovery/holdout protocol to test single layers and multi-layer
combinations. The study quantifies output sensitivity while showing that
apparently favourable search results need not generalise to unseen images.
